\IfFileExists{stacks-project.cls}{%
\documentclass{stacks-project}
}{%
\documentclass{amsart}
}

% For dealing with references we use the comment environment
\usepackage{verbatim}
\newenvironment{reference}{\comment}{\endcomment}
%\newenvironment{reference}{}{}
\newenvironment{slogan}{\comment}{\endcomment}
\newenvironment{history}{\comment}{\endcomment}

% For commutative diagrams we use Xy-pic
\usepackage[all]{xy}

% We use 2cell for 2-commutative diagrams.
\xyoption{2cell}
\UseAllTwocells

% We use multicol for the list of chapters between chapters
\usepackage{multicol}

% This is generally recommended for better output
\usepackage{lmodern}
\usepackage[T1]{fontenc}

% For cross-file-references
\usepackage{xr-hyper}

% Package for hypertext links:
\usepackage{hyperref}

% For any local file, say "hello.tex" you want to link to please
% use \externaldocument[hello-]{hello}
\externaldocument[introduction-]{introduction}
\externaldocument[conventions-]{conventions}
\externaldocument[sets-]{sets}
\externaldocument[categories-]{categories}
\externaldocument[topology-]{topology}
\externaldocument[sheaves-]{sheaves}
\externaldocument[sites-]{sites}
\externaldocument[stacks-]{stacks}
\externaldocument[fields-]{fields}
\externaldocument[algebra-]{algebra}
\externaldocument[brauer-]{brauer}
\externaldocument[homology-]{homology}
\externaldocument[derived-]{derived}
\externaldocument[simplicial-]{simplicial}
\externaldocument[more-algebra-]{more-algebra}
\externaldocument[smoothing-]{smoothing}
\externaldocument[modules-]{modules}
\externaldocument[sites-modules-]{sites-modules}
\externaldocument[injectives-]{injectives}
\externaldocument[cohomology-]{cohomology}
\externaldocument[sites-cohomology-]{sites-cohomology}
\externaldocument[dga-]{dga}
\externaldocument[dpa-]{dpa}
\externaldocument[sdga-]{sdga}
\externaldocument[hypercovering-]{hypercovering}
\externaldocument[schemes-]{schemes}
\externaldocument[constructions-]{constructions}
\externaldocument[properties-]{properties}
\externaldocument[morphisms-]{morphisms}
\externaldocument[coherent-]{coherent}
\externaldocument[divisors-]{divisors}
\externaldocument[limits-]{limits}
\externaldocument[varieties-]{varieties}
\externaldocument[topologies-]{topologies}
\externaldocument[descent-]{descent}
\externaldocument[perfect-]{perfect}
\externaldocument[more-morphisms-]{more-morphisms}
\externaldocument[flat-]{flat}
\externaldocument[groupoids-]{groupoids}
\externaldocument[more-groupoids-]{more-groupoids}
\externaldocument[etale-]{etale}
\externaldocument[chow-]{chow}
\externaldocument[intersection-]{intersection}
\externaldocument[pic-]{pic}
\externaldocument[weil-]{weil}
\externaldocument[adequate-]{adequate}
\externaldocument[dualizing-]{dualizing}
\externaldocument[duality-]{duality}
\externaldocument[discriminant-]{discriminant}
\externaldocument[derham-]{derham}
\externaldocument[local-cohomology-]{local-cohomology}
\externaldocument[algebraization-]{algebraization}
\externaldocument[curves-]{curves}
\externaldocument[resolve-]{resolve}
\externaldocument[models-]{models}
\externaldocument[functors-]{functors}
\externaldocument[equiv-]{equiv}
\externaldocument[pione-]{pione}
\externaldocument[etale-cohomology-]{etale-cohomology}
\externaldocument[proetale-]{proetale}
\externaldocument[relative-cycles-]{relative-cycles}
\externaldocument[more-etale-]{more-etale}
\externaldocument[trace-]{trace}
\externaldocument[crystalline-]{crystalline}
\externaldocument[spaces-]{spaces}
\externaldocument[spaces-properties-]{spaces-properties}
\externaldocument[spaces-morphisms-]{spaces-morphisms}
\externaldocument[decent-spaces-]{decent-spaces}
\externaldocument[spaces-cohomology-]{spaces-cohomology}
\externaldocument[spaces-limits-]{spaces-limits}
\externaldocument[spaces-divisors-]{spaces-divisors}
\externaldocument[spaces-over-fields-]{spaces-over-fields}
\externaldocument[spaces-topologies-]{spaces-topologies}
\externaldocument[spaces-descent-]{spaces-descent}
\externaldocument[spaces-perfect-]{spaces-perfect}
\externaldocument[spaces-more-morphisms-]{spaces-more-morphisms}
\externaldocument[spaces-flat-]{spaces-flat}
\externaldocument[spaces-groupoids-]{spaces-groupoids}
\externaldocument[spaces-more-groupoids-]{spaces-more-groupoids}
\externaldocument[bootstrap-]{bootstrap}
\externaldocument[spaces-pushouts-]{spaces-pushouts}
\externaldocument[spaces-chow-]{spaces-chow}
\externaldocument[groupoids-quotients-]{groupoids-quotients}
\externaldocument[spaces-more-cohomology-]{spaces-more-cohomology}
\externaldocument[spaces-simplicial-]{spaces-simplicial}
\externaldocument[spaces-duality-]{spaces-duality}
\externaldocument[formal-spaces-]{formal-spaces}
\externaldocument[restricted-]{restricted}
\externaldocument[spaces-resolve-]{spaces-resolve}
\externaldocument[formal-defos-]{formal-defos}
\externaldocument[defos-]{defos}
\externaldocument[cotangent-]{cotangent}
\externaldocument[examples-defos-]{examples-defos}
\externaldocument[algebraic-]{algebraic}
\externaldocument[examples-stacks-]{examples-stacks}
\externaldocument[stacks-sheaves-]{stacks-sheaves}
\externaldocument[criteria-]{criteria}
\externaldocument[artin-]{artin}
\externaldocument[quot-]{quot}
\externaldocument[stacks-properties-]{stacks-properties}
\externaldocument[stacks-morphisms-]{stacks-morphisms}
\externaldocument[stacks-limits-]{stacks-limits}
\externaldocument[stacks-cohomology-]{stacks-cohomology}
\externaldocument[stacks-perfect-]{stacks-perfect}
\externaldocument[stacks-introduction-]{stacks-introduction}
\externaldocument[stacks-more-morphisms-]{stacks-more-morphisms}
\externaldocument[stacks-geometry-]{stacks-geometry}
\externaldocument[moduli-]{moduli}
\externaldocument[moduli-curves-]{moduli-curves}
\externaldocument[examples-]{examples}
\externaldocument[exercises-]{exercises}
\externaldocument[guide-]{guide}
\externaldocument[desirables-]{desirables}
\externaldocument[coding-]{coding}
\externaldocument[obsolete-]{obsolete}
\externaldocument[fdl-]{fdl}
\externaldocument[index-]{index}

% Theorem environments.
%
\theoremstyle{plain}
\newtheorem{theorem}[subsection]{Theorem}
\newtheorem{proposition}[subsection]{Proposition}
\newtheorem{lemma}[subsection]{Lemma}

\theoremstyle{definition}
\newtheorem{definition}[subsection]{Definition}
\newtheorem{example}[subsection]{Example}
\newtheorem{exercise}[subsection]{Exercise}
\newtheorem{situation}[subsection]{Situation}

\theoremstyle{remark}
\newtheorem{remark}[subsection]{Remark}
\newtheorem{remarks}[subsection]{Remarks}

\numberwithin{equation}{subsection}

% Macros
%
\def\lim{\mathop{\mathrm{lim}}\nolimits}
\def\colim{\mathop{\mathrm{colim}}\nolimits}
\def\Spec{\mathop{\mathrm{Spec}}}
\def\Hom{\mathop{\mathrm{Hom}}\nolimits}
\def\Ext{\mathop{\mathrm{Ext}}\nolimits}
\def\SheafHom{\mathop{\mathcal{H}\!\mathit{om}}\nolimits}
\def\SheafExt{\mathop{\mathcal{E}\!\mathit{xt}}\nolimits}
\def\Sch{\mathit{Sch}}
\def\Mor{\mathop{\mathrm{Mor}}\nolimits}
\def\Ob{\mathop{\mathrm{Ob}}\nolimits}
\def\Sh{\mathop{\mathit{Sh}}\nolimits}
\def\NL{\mathop{N\!L}\nolimits}
\def\CH{\mathop{\mathrm{CH}}\nolimits}
\def\proetale{{pro\text{-}\acute{e}tale}}
\def\etale{{\acute{e}tale}}
\def\QCoh{\mathit{QCoh}}
\def\Ker{\mathop{\mathrm{Ker}}}
\def\Im{\mathop{\mathrm{Im}}}
\def\Coker{\mathop{\mathrm{Coker}}}
\def\Coim{\mathop{\mathrm{Coim}}}

% Boxtimes
%
\DeclareMathSymbol{\boxtimes}{\mathbin}{AMSa}{"02}

%
% Macros for moduli stacks/spaces
%
\def\QCohstack{\mathcal{QC}\!\mathit{oh}}
\def\Cohstack{\mathcal{C}\!\mathit{oh}}
\def\Spacesstack{\mathcal{S}\!\mathit{paces}}
\def\Quotfunctor{\mathrm{Quot}}
\def\Hilbfunctor{\mathrm{Hilb}}
\def\Curvesstack{\mathcal{C}\!\mathit{urves}}
\def\Polarizedstack{\mathcal{P}\!\mathit{olarized}}
\def\Complexesstack{\mathcal{C}\!\mathit{omplexes}}
% \Pic is the operator that assigns to X its picard group, usage \Pic(X)
% \Picardstack_{X/B} denotes the Picard stack of X over B
% \Picardfunctor_{X/B} denotes the Picard functor of X over B
\def\Pic{\mathop{\mathrm{Pic}}\nolimits}
\def\Picardstack{\mathcal{P}\!\mathit{ic}}
\def\Picardfunctor{\mathrm{Pic}}
\def\Deformationcategory{\mathcal{D}\!\mathit{ef}}
\usepackage{cleveref}
\crefname{originalremark}{original remark}{original remarks}
\Crefname{originalremark}{Original remark}{Original remarks}
\makeatletter
\newlabel{simplicial-remark-relations-cosimplicial}{{016L}{source}{}{sourcepointer-1}{}}
\newlabel{simplicial-remark-relations-cosimplicial@cref}{{[originalremark][0][]016L}{[1][0][]source}{}{}{}}
\makeatother
\crefname{originalsection}{original section}{original sections}
\Crefname{originalsection}{Original section}{Original sections}
\makeatletter
\newlabel{simplicial-section-cosimplicial-object}{{016I}{source}{}{sourcepointer-2}{}}
\newlabel{simplicial-section-cosimplicial-object@cref}{{[originalsection][0][]016I}{[1][0][]source}{}{}{}}
\makeatother
\title{A descent datum gives a cosimplicial module, functorially}
\author{The Stacks Project Authors}
\date{stacks-project:descent.tex:local-snapshot-20261012; GNU FDL 1.2 or later}
\begin{document}
\maketitle
\noindent\textit{Editorial note.} All delivered mathematical text is the author's own, copied from the pinned Stacks Project chapter file descent.tex; no statement or proof step was rewritten, reordered or removed.\par\medskip
\noindent\textit{Editorial note.} The statement carries the author's own footnote about notation for the cosimplicial object; it lies inside the counted statement span and is reproduced verbatim.\par\medskip
\noindent\textit{Editorial note.} The counted spans carry 2 navigation stubs; no bibliography citation occurs in the spans.\par\medskip
\noindent\textit{Editorial note.} The preamble span points at the chapter's own preamble.tex shipped with the frozen source, used verbatim as staged.\par\medskip
\noindent\textit{Allowed standard prerequisites (editorial).}\begin{itemize}
\item Cosimplicial objects and the list of cosimplicial relations. Simplicial, Section simplicial-section-cosimplicial-object and Remark simplicial-remark-relations-cosimplicial, invoked at proof lines 626-627 and 660-662. Supplies the axioms that must be verified.
\item The cocycle condition for a descent datum. Same chapter, the description of descent data for a ring map, used at proof lines 632-654. Supplies the identity that makes the second codegeneracy-coface relation hold.
\item The tensor constructions N\_\{n,n\} attached to a ring map. Same chapter, the notation for the tensor powers of A over R acting on N, used in the statement. Supplies the objects in each cosimplicial degree.
\end{itemize}\medskip
\ifdefined\hypertarget\hypertarget{sourcepointer-1}{}\fi
\noindent\textit{Source reference.} Original remark 016L: sources/stacks-chow/tags.txt, source lines 1548--1548. Pointer outside the counted span, invoked once for context; the load-bearing argument of the counted statement is the author's proof reproduced in full in the proof span.\par\medskip
\ifdefined\hypertarget\hypertarget{sourcepointer-2}{}\fi
\noindent\textit{Source reference.} Original section 016I: sources/stacks-chow/tags.txt, source lines 1545--1545. Pointer outside the counted span, invoked once for context; the load-bearing argument of the counted statement is the author's proof reproduced in full in the proof span.\par\medskip
\begin{lemma}
\label{lemma-descent-datum-cosimplicial}
Let $R \to A$ be a ring map.
Given a descent datum $(N, \varphi)$ we can associate to it a
cosimplicial $(A/R)_\bullet$-module $N_\bullet$\footnote{We should really
write $(N, \varphi)_\bullet$.} by the
rules $N_n = N_{n, n}$ and given $\beta : [n] \to [m]$
setting we define
$$
N_\bullet(\beta) = (\varphi^m_{\beta(n)m}) \circ N_{\beta, n} :
N_{n, n} \longrightarrow N_{m, m}.
$$
This procedure is functorial in the descent datum.
\end{lemma}

\begin{proof}
Here are the first few maps
where $\varphi(n \otimes 1) = \sum \alpha_i \otimes x_i$
$$
\begin{matrix}
\delta^1_0 & : & N & \to & A \otimes N & n & \mapsto & 1 \otimes n \\
\delta^1_1 & : & N & \to & A \otimes N & n & \mapsto &
\sum \alpha_i \otimes x_i\\
\sigma^0_0 & : & A \otimes N & \to & N & a_0 \otimes n & \mapsto & a_0n \\
\delta^2_0 & : & A \otimes N & \to & A \otimes A \otimes N &
a_0 \otimes n & \mapsto & 1 \otimes a_0 \otimes n \\
\delta^2_1 & : & A \otimes N & \to & A \otimes A \otimes N &
a_0 \otimes n & \mapsto & a_0 \otimes 1 \otimes n \\
\delta^2_2 & : & A \otimes N & \to & A \otimes A \otimes N &
a_0 \otimes n & \mapsto & \sum a_0 \otimes \alpha_i \otimes x_i \\
\sigma^1_0 & : & A \otimes A \otimes N & \to & A \otimes N &
a_0 \otimes a_1 \otimes n & \mapsto & a_0a_1 \otimes n \\
\sigma^1_1 & : & A \otimes A \otimes N & \to & A \otimes N &
a_0 \otimes a_1 \otimes n & \mapsto & a_0 \otimes a_1n
\end{matrix}
$$
with notation as in
Simplicial, Section \ref{simplicial-section-cosimplicial-object}.
We first verify the two properties $\sigma^0_0 \circ \delta^1_0 = \text{id}$
and $\sigma^0_0 \circ \delta^1_1 = \text{id}$.
The first one, $\sigma^0_0 \circ \delta^1_0 = \text{id}$, is clear from
the explicit description of the morphisms above.
To prove the second relation we have to use the cocycle condition
(because it does not hold for an arbitrary isomorphism
$\varphi : N \otimes_R A \to A \otimes_R N$). Write
$p = \sigma^0_0 \circ \delta^1_1 : N \to N$. By the description of the
maps above we deduce that $p$ is also equal to
$$
p = \varphi \otimes \text{id} :
N = (N \otimes_R A) \otimes_{(A \otimes_R A)} A
\longrightarrow
(A \otimes_R N) \otimes_{(A \otimes_R A)} A = N
$$
Since $\varphi$ is an isomorphism we see that $p$ is an isomorphism.
Write $\varphi(n \otimes 1) = \sum \alpha_i \otimes x_i$ for certain
$\alpha_i \in A$ and $x_i \in N$. Then $p(n) = \sum \alpha_ix_i$.
Next, write
$\varphi(x_i \otimes 1) = \sum \alpha_{ij} \otimes y_j$ for
certain $\alpha_{ij} \in A$ and $y_j \in N$. Then the cocycle condition
says that
$$
\sum \alpha_i \otimes \alpha_{ij} \otimes y_j
=
\sum \alpha_i \otimes 1 \otimes x_i.
$$
This means that $p(n) = \sum \alpha_ix_i = \sum \alpha_i\alpha_{ij}y_j =
\sum \alpha_i p(x_i) = p(p(n))$. Thus $p$ is a projector, and since it is
an isomorphism it is the identity.

\medskip\noindent
To prove fully that $N_\bullet$ is a cosimplicial module we have to check
all 5 types of relations of
Simplicial, Remark \ref{simplicial-remark-relations-cosimplicial}.
The relations on composing $\sigma$'s are obvious.
The relations on composing $\delta$'s come down to the
cocycle condition for $\varphi$.
In exactly the same way as above one checks the relations
$\sigma_j \circ \delta_j = \sigma_j \circ \delta_{j + 1} = \text{id}$.
Finally, the other relations on compositions of $\delta$'s and $\sigma$'s
hold for any $\varphi$ whatsoever.
\end{proof}

\end{document}
